\section{Discussion and Limitations}
\label{sec:discussion}

\paragraph{What the evidence supports.}
The measurement results are the firm part of this paper. They are counterfactual, repeated on disjoint
game sets and on a trained policy, and come with intervals. Inside an agent's reasoning the decision is
made before the command is written; uncertainty signals (entropy, NLL) do not locate it; the commitment
probe does. GACA's criticality score is dominated by within-state path noise. The training results are
weaker. Probe-weighted credit had the highest final success, and the variants that changed how it is
applied did worse, two of them significantly. With one seed per arm and about 6\% of the published
training budget, no difference between the main methods is statistically resolved.

\paragraph{Why a better measurement need not train better.}
The probe changes where credit lands, not whether it is right. At 5--15\% training success most step
advantages rest on one or two successful rollouts in a group of eight, so their sign is often wrong;
concentrating a wrong-signed advantage on the reasoning that produced the decision does more damage than
spreading it. Uniform credit is, in this sense, a hedge against a noisy step estimator, and the positive-
only variant shows that asymmetric hedging is worse than none. We expect probe weighting to pay off more
where step advantages are reliable: at higher success rates, with more rollouts per group, or combined
with an estimator that scales concentration by its own confidence, for instance the number of successful
rollouts in the group or whether the anchor group carries a real return contrast.

\paragraph{Limitations.}
\emph{Budget and seeds}: about 6\% of GACA's trajectories and one seed per method, so method
differences of a few points are not interpretable on their own; the paired tests are the right lens,
and they are mostly inconclusive. \emph{Training seed}: the probe arm drew its training games in a
different order than the controls, with identical evaluation games. \emph{Benchmark}: ALFWorld only;
WebShop, the paper's second benchmark, needs a product-search server we did not stand up.
\emph{Oracle}: the flip test measures relevance for the next command only, not for the episode's
outcome, and uses 8 resamples per cut; importance is a difference of two noisy rates. \emph{Probe}: as a
likelihood ratio under teacher forcing, the probe can credit a segment that makes the eventual command
\emph{expressible} rather than one that decides it. \emph{Cost}: the probe adds about 40\% to an iteration
in our setup; computing it from a prefix-shared batch would reduce this.

\paragraph{Outlook.}
The cheapest next steps follow directly: multiple seeds of the probe and GACA at this budget, the
paper's 16-task batch, and probe weights scaled by a confidence in the step advantage. More broadly, our
results suggest a separation of concerns for critic-free agent RL: the step estimator decides
\emph{how much} credit a step earns, and a counterfactual probe decides \emph{which tokens} earned it.
The second is now cheap and measurable; the first is the bottleneck.

\section{Conclusion}
We asked where, inside an agent's reasoning, the decision is made. A counterfactual flip test answers
that it is made in the reasoning, mostly not where the model is uncertain, and that a teacher-forced
commitment probe, the measurement behind CIR, finds it reliably. We turned the probe into a
mean-preserving redistribution of GACA's step advantage over tokens. In a budget-matched, game-paired
study on ALFWorld it achieved the best final success of the methods we trained, but the margins are
within single-seed noise, and ways of applying it that seemed safer did worse. Measuring where credit
belongs has become easy; making the credit itself reliable is the open problem.
